\documentclass[UTF8,a4paper,11pt]{ctexart}

\usepackage[margin=2.45cm,headheight=14pt]{geometry}
\usepackage{amsmath}
\usepackage{array}
\usepackage{booktabs}
\usepackage{enumitem}
\usepackage{fancyhdr}
\usepackage{hyperref}
\usepackage{longtable}
\usepackage{tabularx}
\usepackage{tikz}
\usepackage{xcolor}
\usepackage{xurl}

\usetikzlibrary{arrows.meta,positioning,fit}

\definecolor{ReportBlue}{HTML}{24527A}
\definecolor{ReportGreen}{HTML}{34745B}
\definecolor{ReportGold}{HTML}{A66A1F}
\definecolor{ReportRed}{HTML}{9A3E3E}
\definecolor{ReportGray}{HTML}{5A6470}
\definecolor{ReportLight}{HTML}{F3F6F8}

\hypersetup{
  colorlinks=true,
  linkcolor=ReportBlue,
  urlcolor=ReportBlue,
  citecolor=ReportGreen,
  pdftitle={从 AI 算力到超节点：昇腾平台上的 TileXR EP Combine 算子开发实践},
  pdfauthor={实习月度报告}
}

\setlength{\parindent}{2em}
\setlength{\parskip}{0.25em}
\setlist{itemsep=0.25em,topsep=0.35em}
\renewcommand{\arraystretch}{1.28}

\pagestyle{fancy}
\fancyhf{}
\fancyhead[L]{TileXR EP Combine 算子开发实践}
\fancyhead[R]{2026 年 7 月实习报告}
\fancyfoot[C]{\thepage}

\newcommand{\code}[1]{\texttt{#1}}
\newcommand{\term}[1]{\textbf{#1}}
\newcommand{\placeholder}[1]{\underline{\hspace{#1}}}

\title{\vspace{-1.5cm}\textbf{从 AI 算力到超节点：\\[0.45em]
昇腾平台上的 TileXR EP Combine 算子开发实践}}
\author{
  \begin{tabular}{rl}
    实习生： & \placeholder{4.5cm} \\
    学校/专业： & \placeholder{4.5cm} \\
    实习部门： & \placeholder{4.5cm} \\
    指导老师： & \placeholder{4.5cm}
  \end{tabular}
}
\date{2026 年 7 月}

\begin{document}

\maketitle
\thispagestyle{empty}
\vfill

\begin{center}
\colorbox{ReportLight}{
  \parbox{0.88\textwidth}{
    \textbf{报告关键词：}AI 计算、超节点、昇腾、Ascend C、Unified Buffer、
    MoE、专家并行、Combine、UDMA/URMA、正确性验证、性能分析
  }
}
\end{center}

\vfill
\newpage

\begin{abstract}
本月实习围绕昇腾平台上的高性能通信算子开发展开。学习过程并不是从一个复杂算子的
源代码直接开始，而是沿着“理解硬件架构、掌握 Ascend C 编程模型、认识片上数据搬运
与卡间通信、完成真实算子实践”的路线逐步深入。最终，我在 TileXR 项目中完成了
专家并行（Expert Parallelism，EP）场景下 Combine 算子的开发与实验，建立了从代码
修改、独立编译、正确性验证到性能分析的完整工作流程。

为便于非专业读者理解，本文首先从当前 AI 计算面临的算力、存储和通信问题出发，说明
为什么多芯片协同和“超节点”逐渐成为大模型系统的重要基础设施；随后以通俗方式介绍
昇腾 AI Core、Ascend C、Unified Buffer（UB）以及 UDMA/URMA 通信；最后结合
TileXR Combine 算子，说明一个分布式 AI 算子如何把计算、数据搬运、同步协议和性能
测量组织成可验证、可复现的工程系统。本文的重点不仅是“完成了什么”，也包括“如何
确认它是正确的、如何判断它是否真的更快”。
\end{abstract}

\tableofcontents
\newpage

\section{本月工作概览}

本月的学习和实践可以概括为四个递进阶段，如表~\ref{tab:overview} 所示。前两个阶段
建立对平台的认识，第三阶段把知识连接为一个完整的数据通路，第四阶段则在真实项目中
完成闭环。四个阶段之间并不是相互独立的：没有对硬件存储层次的理解，就很难解释
数据为什么要经过 UB；没有对同步和通信的理解，就无法保证多卡 Combine 在重复运行
时仍然正确；没有严格的测量方法，也无法判断一次代码改动究竟带来了收益还是噪声。

\begin{table}[htbp]
  \centering
  \caption{本月学习与实践主线}
  \label{tab:overview}
  \begin{tabularx}{\textwidth}{>{\bfseries}p{0.15\textwidth}X X}
    \toprule
    阶段 & 主要内容 & 阶段产出 \\
    \midrule
    平台认识 & 学习昇腾硬件架构、AI Core 和存储层次 & 能够解释 GM、UB、计算流水线之间的数据关系 \\
    算子入门 & 学习 Ascend C 的 Host/Kernel 分工、Tiling、数据搬运与同步 & 能够阅读并修改算子代码，理解编译和运行链路 \\
    通信学习 & 学习片上流水线同步、控制标志以及 UDMA/URMA 卡间传输 & 能够从生产者、消费者和可见性角度分析通信协议 \\
    项目实践 & 开发 TileXR EP Combine，并开展正确性和性能实验 & 形成代码修改、独立编译、正确性验证、性能分析闭环 \\
    \bottomrule
  \end{tabularx}
\end{table}

从结果上看，本月最大的收获并不只是增加了一个算子实现，而是形成了一个可重复使用的
工程方法：先定义问题和计时边界，再修改代码；先验证正确性，再比较性能；对每次实验
保存配置、日志和原始数据，避免只保留一个看起来更好的数字。

\section{从当前 AI 计算现状说起}

\subsection{AI 计算为什么越来越依赖系统能力}

现代 AI 模型的能力提升通常伴随着更大的参数量、更长的上下文和更多的训练、推理数据。
从表面上看，这是“需要更多算力”；但在真实系统中，瓶颈往往不只来自乘加运算，还来自
以下三个方面：

\begin{enumerate}
  \item \term{容量问题：}单颗芯片的显存有限，大模型的参数、激活和中间结果无法全部容纳；
  \item \term{带宽问题：}计算单元速度很快，如果数据不能及时到达，计算单元就会等待；
  \item \term{协同问题：}模型分布在多颗芯片后，芯片之间需要交换数据并保持一致的执行节奏。
\end{enumerate}

可以把一颗 AI 芯片理解为一家加工能力很强的工厂。工厂中的机器再快，如果仓库取料慢、
运输车辆堵塞，最终产量仍然受物流限制。大模型系统也是如此：计算、存储和通信必须共同
设计，单独提高其中一项并不能保证端到端性能同步提高。

\subsection{MoE：只激活一部分专家}

混合专家模型（Mixture of Experts，MoE）是一种扩展模型容量的重要方法。模型中可以
设置多个“专家”网络，但每个 token 通常只被路由到其中少数几个专家。这样，模型可以
拥有较大的总参数规模，同时把一次计算的活跃计算量控制在合理范围内\cite{shazeer2017moe,fedus2021switch}。

MoE 引入了新的系统问题：专家可能分布在不同芯片上，因此 token 要先被发送到对应专家，
专家计算完成后再把结果送回原位置。前一个过程通常称为 \term{Dispatch}，后一个过程称为
\term{Combine}。当模型规模扩大到多卡甚至多机时，这两个步骤会形成密集的多对多通信，
通信效率便直接影响模型整体效率。

\subsection{为什么会引出超节点}

传统集群通过网络连接多台服务器，扩展灵活，但跨设备通信的带宽和时延与芯片内部差距
很大。为了让多颗加速芯片更像一个整体，业界开始构建更紧密的计算单元：在一个受控范围
内提供高带宽、低时延互连，并配合统一的通信软件和任务调度。这类系统通常被称为
\term{超节点}。

超节点不是简单地“把更多卡装在一起”。它强调的是计算、内存、互连和软件栈的协同：

\begin{itemize}
  \item 对应用而言，多颗芯片应能高效协作，而不是彼此孤立；
  \item 对通信而言，要减少不必要的数据复制和 Host 介入；
  \item 对算子而言，要把计算、搬运和同步重叠起来，减少等待；
  \item 对工程而言，要能定位最慢的芯片、最慢的通信路径和最慢的流水线阶段。
\end{itemize}

因此，从 AI 计算现状到超节点，再到本月的 Combine 实践，逻辑关系可以概括为：

\begin{figure}[htbp]
  \centering
  \begin{tikzpicture}[
    node distance=0.7cm,
    box/.style={draw,rounded corners=2pt,minimum height=1.05cm,text width=3.1cm,
      align=center,fill=ReportLight},
    arrow/.style={-{Latex[length=2.4mm]},thick,draw=ReportBlue}
  ]
    \node[box] (model) {模型规模与数据量增长};
    \node[box,right=of model] (multi) {单芯片容量与带宽受限};
    \node[box,right=of multi] (super) {多芯片紧密协同\\超节点};
    \node[box,right=of super] (op) {通信计算融合\\Combine 算子};
    \draw[arrow] (model) -- (multi);
    \draw[arrow] (multi) -- (super);
    \draw[arrow] (super) -- (op);
  \end{tikzpicture}
  \caption{从 AI 计算需求到 Combine 算子实践的逻辑链条}
  \label{fig:motivation}
\end{figure}

\section{昇腾架构与 Ascend C 算子开发}

\subsection{从“计算芯片”到“数据通路”}

昇腾 NPU 面向 AI 计算提供专用计算核心。对本次实习最重要的认识是：开发算子时不能只
关注公式，还要关注数据位于哪里、由谁搬运、何时可见。一个简化的数据层次如下：

\begin{figure}[htbp]
  \centering
  \begin{tikzpicture}[
    level/.style={draw,minimum height=1.0cm,align=center,rounded corners=2pt},
    arrow/.style={<->,>=Latex,thick,draw=ReportGray}
  ]
    \node[level,fill=ReportLight,text width=3.2cm] (gm) {Global Memory（GM）\\容量较大，访问较慢};
    \node[level,fill=ReportLight,text width=3.2cm,right=1.1cm of gm] (ub) {Unified Buffer（UB）\\片上暂存，容量较小};
    \node[level,fill=ReportLight,text width=3.2cm,right=1.1cm of ub] (calc) {Vector/Scalar 等流水线\\执行计算与控制};
    \draw[arrow] (gm) -- node[above]{MTE 搬运} (ub);
    \draw[arrow] (ub) -- node[above]{读写数据} (calc);
  \end{tikzpicture}
  \caption{面向算子开发的简化数据通路}
  \label{fig:memory}
\end{figure}

其中，GM 可以类比为工厂的大仓库，UB 可以类比为生产线旁边的小料架。大仓库容量大，
但每次取料成本较高；小料架离计算单元近，但容量有限。高性能算子需要把数据分块搬入
UB，在计算的同时准备下一块数据，并在合适时机把结果写回 GM。

这里需要特别澄清：\term{UB（Unified Buffer）本身是片上缓冲存储，并不是网络通信
协议}。本月学习的“UB 相关通信”更准确地说，是围绕 GM 与 UB 的数据搬运、不同流水线
之间的事件同步，以及在此基础上延伸出的设备间通信与可见性协议。

\subsection{Ascend C 的编程视角}

Ascend C 为算子开发提供接近 C++ 的编程方式，同时暴露 AI Core 的数据搬运和计算能力。
一个完整算子通常包含两类工作：

\begin{itemize}
  \item \term{Host 侧：}检查输入、计算 Tiling 参数、准备 Workspace、选择 Kernel 配置并发起执行；
  \item \term{Kernel 侧：}在 AI Core 上完成数据搬运、向量计算、核间同步和结果写回。
\end{itemize}

本月在阅读和修改代码时，我逐步理解了
\code{DataCopy}/\code{DataCopyPad}、\code{Cast}、向量计算、事件标志、流水线屏障和双缓冲
等概念。与普通 CPU 程序相比，算子开发更强调显式的数据布局和执行顺序。例如，数据从
GM 搬到 UB 后，必须等待对应搬运事件完成，后续 Scalar 或 Vector 流水线才能安全读取；
输出写回 GM 前，也必须保证向量计算已经结束。

\subsection{流水线与双缓冲}

如果每一轮都严格按照“搬入、等待、计算、等待、搬出”的顺序串行执行，硬件中会出现
大量空闲时间。双缓冲的思路是准备两块 UB 空间：计算第 0 块数据时，可以同时搬运第 1 块；
下一轮交换两块空间的角色。其目标不是减少计算量，而是把搬运时间隐藏在计算时间中。

不过，双缓冲也带来了新的正确性要求：一块缓冲区只有在上一位消费者完全使用结束后
才能复用。因此，“并行得更多”必须与“同步得正确”同时满足，这一认识后来直接应用到
Combine 的发送、接收和窗口复用设计中。

\section{从片上同步到卡间通信}

\subsection{通信不仅是把数据从 A 复制到 B}

在单线程程序中，写完数据再读取通常是自然发生的；在多核、多卡环境中，生产者和消费者
可能独立运行。一个完整的通信过程至少要回答四个问题：

\begin{enumerate}
  \item 数据写到了什么地址；
  \item 消费者如何知道数据已经写完；
  \item “看到完成标志”是否意味着也能看到完整数据；
  \item 这块内存何时可以被下一轮安全复用。
\end{enumerate}

这也是为什么 Combine 不能只调用一次复制接口。数据传输之外，还需要 Ready 标志、
内存屏障、完成确认和轮次编号共同构成通信协议。

\subsection{UDMA/URMA 的作用}

TileXR 在目标平台上使用 UDMA/URMA 能力，让设备侧能够针对已注册的远端内存直接提交
数据传输。可以把它理解为：AI Core 不必把数据先交回 CPU，再由 CPU 转发，而是能够
准备传输描述、敲响 doorbell，并等待队列完成。这样既减少 Host 参与，也为计算与通信
重叠提供了条件。

非阻塞提交只表示“传输任务已经交给硬件”，并不等于“远端已经收到完整数据”。因此，
实现中还要通过完成队列和 Quiet 等机制确认传输完成，再发布本轮结束状态。这个区别是
本月理解设备通信语义的重要一步。

\subsection{三层 Ready 协议}

Combine 实践中使用的同步可以用三层关系解释：

\begin{longtable}{>{\bfseries}p{0.18\textwidth}p{0.27\textwidth}p{0.46\textwidth}}
  \caption{Combine 中三类 Ready/控制机制}\label{tab:ready}\\
  \toprule
  机制 & 连接的角色 & 解决的问题 \\
  \midrule
  \endfirsthead
  \toprule
  机制 & 连接的角色 & 解决的问题 \\
  \midrule
  \endhead
  TxReady & 本地 Pack 与本地 Send & Pack 完整写好一条待发送数据后，Send 才能读取并提交传输 \\
  DataAsFlag & 远端 Send/UDMA 与 Receive & Receive 看到数据块尾部标志后，才读取对应 payload \\
  Control Line & 多核与多 rank & 协调启动、发送完成、接收完成和窗口释放等全局状态 \\
  \bottomrule
\end{longtable}

每个控制值都带有当前轮次的 generation（在实现中由 \code{magic} 表示）。这样，即使内存
中还保留着上一轮的标志，新一轮也不会把旧状态误认为当前状态。这个设计体现了通信协议
的核心：标志不仅要表示“完成”，还必须说明“是哪一轮的完成”。

\section{TileXR EP Combine 算子实践}

\subsection{Combine 要完成什么工作}

在专家并行的 MoE 模型中，Dispatch 根据路由结果把 token 发往不同专家所在的 rank。
专家计算后，每个原始 token 会得到若干个专家输出。Combine 的任务是把这些输出送回
token 所属的 rank，按权重还原并累加：

\begin{equation}
  y_{t}=\sum_{k=1}^{K} w_{t,k}\cdot
  \operatorname{dequant}\!\left(e_{t,k}\right),
\end{equation}

其中 $t$ 表示 token，$K$ 表示被选中的专家数量，$e_{t,k}$ 是第 $k$ 个专家输出，
$w_{t,k}$ 是对应的路由权重。公式本身并不复杂，困难在于这些专家输出可能来自不同 rank，
到达时间也不完全一致。

\begin{figure}[htbp]
  \centering
  \begin{tikzpicture}[
    node distance=0.65cm,
    box/.style={draw,rounded corners=2pt,minimum height=0.95cm,text width=2.3cm,
      align=center,fill=ReportLight},
    arrow/.style={-{Latex[length=2.3mm]},thick,draw=ReportBlue}
  ]
    \node[box] (token) {原始 token};
    \node[box,right=of token] (dispatch) {Dispatch\\发送到专家};
    \node[box,right=of dispatch] (expert) {多 rank\\专家计算};
    \node[box,right=of expert] (combine) {Combine\\回传与加权};
    \node[box,right=of combine] (output) {最终 token 输出};
    \draw[arrow] (token) -- (dispatch);
    \draw[arrow] (dispatch) -- (expert);
    \draw[arrow] (expert) -- (combine);
    \draw[arrow] (combine) -- (output);
  \end{tikzpicture}
  \caption{MoE 中 Dispatch、专家计算与 Combine 的关系}
  \label{fig:moe}
\end{figure}

\subsection{算子的主要数据流}

本次实现把 64 个逻辑 AIV 分成 Pack/Receive 和 Send 两类角色。前者负责将专家结果量化
打包，并在远端数据到达后完成反量化、TopK 加权累加和输出；后者负责扫描本地已就绪的
route，提交 UDMA 传输并参与轮次控制。最终交付配置采用 42 个 Pack/Receive 核与 22 个
Send 核。这个划分不是理论上预先确定的常数，而是通过不同配置的双负载实验比较得到。

\begin{figure}[htbp]
  \centering
  \begin{tikzpicture}[
    node distance=0.55cm,
    stage/.style={draw,rounded corners=2pt,minimum height=1.0cm,text width=2.15cm,
      align=center,fill=ReportLight},
    signal/.style={draw,dashed,rounded corners=2pt,minimum height=0.75cm,text width=1.9cm,
      align=center,fill=white},
    arrow/.style={-{Latex[length=2.2mm]},thick,draw=ReportBlue},
    ctrl/.style={-{Latex[length=2.0mm]},dashed,draw=ReportGold}
  ]
    \node[stage] (pack) {Pack\\量化与打包};
    \node[signal,right=of pack] (tx) {TxReady};
    \node[stage,right=of tx] (send) {Send\\UDMA 提交};
    \node[signal,right=of send] (flag) {DataAsFlag};
    \node[stage,right=of flag] (recv) {Receive\\反量化与累加};
    \draw[arrow] (pack) -- (tx);
    \draw[arrow] (tx) -- (send);
    \draw[arrow] (send) -- (flag);
    \draw[arrow] (flag) -- (recv);
    \node[signal,below=0.85cm of send] (round) {Round Control};
    \draw[ctrl] (send) -- (round);
    \draw[ctrl] (recv) -- (round);
  \end{tikzpicture}
  \caption{Combine 的主要数据与控制路径}
  \label{fig:combine-flow}
\end{figure}

这个数据流包含以下关键步骤：

\begin{enumerate}
  \item \term{Pack：}读取专家输出，按 route 量化并组织为适合通信的连续数据块；
  \item \term{本地发布：}数据完整写回后发布 TxReady，通知 Send 核可以消费；
  \item \term{Send：}读取 metadata，区分本地复制和远端发送，向 UDMA 队列提交任务；
  \item \term{Receive：}轮询远端数据块的 DataAsFlag，确认就绪后搬入 UB；
  \item \term{恢复与聚合：}反量化专家输出，乘以 TopK 权重，并使用 FP32 累加；
  \item \term{轮次收尾：}确认发送、接收和输出全部完成，再发布窗口可复用状态。
\end{enumerate}

\subsection{实现中的重点问题}

\subsubsection{正确的数据可见性顺序}

实现必须保证“先有数据，后有标志”。如果生产者提前发布 Ready，消费者可能读到尚未写完
的数据。反过来，消费者看到 Ready 后也要建立 acquire 顺序，再开始读取 payload。本次
实践将这类要求整理为清晰的不变量，例如：

\begin{center}
  \fbox{\parbox{0.9\textwidth}{
    TxData 完整可见 $\rightarrow$ TxReady 可见 $\rightarrow$ Send 读取；\\
    Rx payload 完整可见 $\rightarrow$ DataAsFlag 可见 $\rightarrow$ Receive 读取；\\
    全部发送完成且全部接收完成 $\rightarrow$ 发布 RxWindow 可复用。
  }}
\end{center}

\subsubsection{Ping-pong 窗口的安全复用}

接收区使用两个窗口按轮次奇偶交替，从而避免相邻轮次立即覆盖同一批地址。但两个窗口
并不意味着可以取消同步：当同一窗口在两轮之后再次被使用时，仍要确认上一位消费者已经
完成并清理标志。这个问题让我认识到，缓冲区数量和内存所有权协议是两个不同维度。

\subsubsection{核角色与负载平衡}

增加 Send 核可以提高同时提交传输的能力，但也会减少执行 Pack/Receive 的核心数量；
反之，Pack/Receive 核过多又可能使发送侧成为瓶颈。因此需要把核角色比例视为端到端配置，
而不是孤立优化某一阶段。最终使用不同 batch size 共同选择配置，避免只对单一负载过拟合。

\section{建立完整的实验流程}

\subsection{代码修改：先明确改动假设}

每次修改前先写清楚它希望解决的具体问题，例如“减少重复 Ready 读取”“提前发布某个状态”
或“调整 Pack/Send 核比例”。一个改动尽量只对应一个主要假设，便于通过 A/B 实验判断效果。
如果一次同时改变多个条件，即使结果变快，也很难知道真正有效的是哪一项。

代码层面同时关注 Host 与 Kernel：Host 侧负责参数、Workspace 布局和启动配置，Kernel 侧
负责设备上的数据路径。两边对偏移、大小和角色划分的理解必须一致，否则错误可能表现为
越界、数据错位，甚至多卡等待无法结束。

\subsection{独立编译：确认测试的就是当前代码}

编译并不只是把源代码变成二进制，还决定了内联、指令布局和优化级别。实践中显式检查
BiSheng 编译命令和 driver 展开的参数，并分别比较不同优化级别。最终交付选择 O2：它保留
主要性能收益，同时通过正确性和双负载性能验证。

为了避免“修改了源代码但运行的仍是旧产物”，实验记录中保存源码、Kernel 和动态库的
哈希，并将构建配置与运行结果放在同一个可追溯目录中。这使编译产物成为实验的一部分，
而不是被忽略的中间文件。

\subsection{正确性验证：不仅比较一次输出}

分布式算子容易出现只在特定轮次、特定 rank 或特定输入规模触发的问题。因此正确性验证
不能只做一次数值比较，而应覆盖表~\ref{tab:correctness} 中的多个维度。

\begin{table}[htbp]
  \centering
  \caption{Combine 正确性验证维度}
  \label{tab:correctness}
  \begin{tabularx}{\textwidth}{>{\bfseries}p{0.23\textwidth}X}
    \toprule
    验证维度 & 目的 \\
    \midrule
    参考结果对比 & 确认量化、反量化、TopK 加权和输出位置正确 \\
    多 batch size & 检查边界、分块和不同 token 分配是否正确 \\
    多 rank 联合运行 & 检查远端寻址、路由信息和全局同步 \\
    多轮重复运行 & 检查 generation、Ping-pong 窗口和旧标志残留 \\
    非法 metadata/header & 检查错误能被记录，同时协议可以安全收尾 \\
    优化级别切换 & 暴露可能被低优化级别掩盖的生命周期或别名问题 \\
    \bottomrule
  \end{tabularx}
\end{table}

正确性优先于性能。如果一个方案只能通过放松同步或减少必要检查来获得更低时间，它并不
构成可接受的优化。性能实验只在正确性 smoke test 通过后进入正式统计。

\subsection{性能分析：先定义“时间”是什么}

性能分析中最容易犯的错误，是把不同边界测得的时间直接比较。本次实践区分了三类指标：

\begin{itemize}
  \item \term{Host 时间：}包含发起、调度以及可能的同步开销，适合观察应用侧体验；
  \item \term{Strict Kernel 时间：}尽量少插桩，表示生产 Kernel 的关键路径；
  \item \term{分阶段 Profile：}记录 Pack、Send、Receive 和等待阶段，用于解释瓶颈，但插桩本身可能改变性能。
\end{itemize}

正式配置选择使用 profiling-off 的 Strict Kernel 指标。聚合顺序为：每个 rank 内取 64 个
AIV 的最大值，每次 launch 再取所有 rank 的最大值，最后对多次有效 launch 取中位数。
这样统计的是整个分布式请求的关键路径，因为最终完成时间由最慢的 core 和最慢的 rank
决定，而不是由平均值决定。

\begin{figure}[htbp]
  \centering
  \begin{tikzpicture}[
    node distance=0.7cm,
    box/.style={draw,rounded corners=2pt,minimum height=0.95cm,text width=3.0cm,
      align=center,fill=ReportLight},
    arrow/.style={-{Latex[length=2.3mm]},thick,draw=ReportGreen}
  ]
    \node[box] (change) {提出假设并修改代码};
    \node[box,right=of change] (build) {独立编译并记录产物};
    \node[box,right=of build] (correct) {正确性与多轮验证};
    \node[box,below=0.8cm of correct] (profile) {Strict 测量与阶段归因};
    \node[box,left=of profile] (compare) {A/B 对比与结果复核};
    \node[box,left=of compare] (archive) {保存配置、日志与原始数据};
    \draw[arrow] (change) -- (build);
    \draw[arrow] (build) -- (correct);
    \draw[arrow] (correct) -- (profile);
    \draw[arrow] (profile) -- (compare);
    \draw[arrow] (compare) -- (archive);
    \draw[arrow] (archive) -- (change);
  \end{tikzpicture}
  \caption{本月建立的算子实验闭环}
  \label{fig:workflow}
\end{figure}

\section{实验结果与分析}

\subsection{固定环境下的代表性结果}

表~\ref{tab:performance} 给出项目交付记录中的代表性结果。测试固定在八卡、隐藏维度
$H=5120$、TopK$=6$ 的负载上，并同时比较 BS32 和 BS128。表中时间保留 raw cycles，
显示微秒值沿用项目测量约定。这里的数字用于比较同一环境下的配置，不用于推导不同硬件
或不同模型上的绝对结论。

\begin{table}[htbp]
  \centering
  \caption{八卡 Strict Kernel 代表性结果（100 次有效 launch 的中位数）}
  \label{tab:performance}
  \begin{tabular}{lrrr}
    \toprule
    核角色配置 & Batch Size & 中位数/cycles & 显示值/$\mu$s \\
    \midrule
    P42/S22 & 32  & 43,130  & 43.130 \\
    P36/S28 & 32  & 46,716  & 46.716 \\
    P42/S22 & 128 & 94,889  & 94.889 \\
    P36/S28 & 128 & 115,397 & 115.397 \\
    \bottomrule
  \end{tabular}
\end{table}

其中 P 表示 Pack/Receive 核数量，S 表示 Send 核数量。P42/S22 相对 P36/S28 在 BS32
上改善约 $7.68\%$，在 BS128 上改善约 $17.77\%$。结果说明：继续增加 Send 核并没有
稳定降低端到端时间，反而会侵占 Pack/Receive 资源。较大的 BS128 对这种资源失衡更敏感，
因此双负载共同选择比只看 BS32 更可靠。

\subsection{如何避免“测得越详细，看起来越慢”}

早期分阶段 Profile 在很多位置记录时间和计数，插桩会增加指令、同步和分支，并改变代码
布局，因此测得的绝对时间明显高于生产 Kernel。为验证这一点，实践中增加了低侵入的
total-only 测量，只记录总时间，并使用与 Strict 相同的边界。在一个代表性配置上，
matched strict10 为 41,703.5 cycles，total-only10 为 41,126 cycles，差异约为
$-1.38\%$。这说明旧 Profile 的大幅膨胀主要来自观测方式，而不是生产 Kernel 的真实成本。

由此形成了明确分工：Strict 用于选择最终配置，分阶段 Profile 用于解释时间花在哪里。
某个阶段单独变快，只能作为线索，不能替代端到端结果和正确性验证。

\subsection{未采用方案同样有价值}

研发过程中也尝试了提前发布 TxReady、批量 Ready 读取、Vector 比较、不同核比例以及不同
编译优化级别等方向。部分方案在某个 batch size 或某个局部阶段上有收益，但没有通过
双负载、重复运行或综合稳定性验证，因此没有进入最终配置。

这些“未采用”结果不是无效工作。它们帮助缩小了瓶颈范围，也证明性能优化不能依赖单次
最好值。保留失败配置和原始记录，使后续开发者能够知道哪些方向已经验证过、为什么没有
合入，避免重复试错。

\section{本月遇到的主要困难与解决方法}

\begin{longtable}{>{\bfseries}p{0.22\textwidth}p{0.34\textwidth}p{0.34\textwidth}}
  \caption{主要困难、分析方法与收获}\label{tab:challenges}\\
  \toprule
  困难 & 处理方法 & 形成的认识 \\
  \midrule
  \endfirsthead
  \toprule
  困难 & 处理方法 & 形成的认识 \\
  \midrule
  \endhead
  概念跨度大 & 从存储层次和最小数据通路入手，再阅读完整 Kernel & 先建立系统地图，再进入局部代码效率更高 \\
  多核多卡错误难复现 & 使用 generation、重复轮次、错误状态和超时信息定位 & 分布式正确性依赖明确的状态机，而不仅是数值输出 \\
  Profile 与 Strict 差距大 & 对齐计时边界，增加 total-only 交叉验证 & 测量工具会影响被测对象，必须控制观测开销 \\
  局部优化未带来总收益 & 同时观察最慢 core、最慢 rank 和双负载结果 & 端到端性能由关键路径决定，不等于局部阶段之和 \\
  编译配置影响结果 & 保存实际编译命令、优化级别与产物哈希 & 编译器和代码生成也是系统设计的一部分 \\
  实验数量多、容易混淆 & 为每个 variant 固化配置并隔离结果目录 & 可追溯性是性能结论可信度的基础 \\
  \bottomrule
\end{longtable}

其中最重要的转变，是从“看代码是否合理”转向“用证据确认系统行为”。例如，一个 Ready
标志从代码顺序上看似乎写在数据之后，但仍要结合流水线事件和内存屏障判断硬件可见性；
一个优化从局部计时看更快，也仍要用端到端 Strict 指标确认其是否缩短关键路径。

\section{阶段总结与能力提升}

\subsection{技术能力}

经过本月实践，我对昇腾算子开发形成了从硬件到应用的分层认识：

\begin{itemize}
  \item 能够从 GM、UB 和执行流水线角度分析数据路径；
  \item 能够区分 Host 侧配置与 Kernel 侧执行职责；
  \item 能够理解多核生产者--消费者同步、卡间传输和窗口复用；
  \item 能够阅读 EP Combine 的 Pack、Send、Receive 和控制流程；
  \item 能够针对真实改动完成编译、正确性验证、性能采集和结果解释。
\end{itemize}

\subsection{工程能力}

相比具体 API，本月更具长期价值的是工程习惯的建立：

\begin{enumerate}
  \item \term{问题可拆分：}把复杂算子拆成数据布局、计算、通信、同步和测量五类问题；
  \item \term{实验可复现：}固定代码、编译参数、运行参数和统计口径；
  \item \term{结论可审计：}保留 raw cycles、逐 rank 结果和失败实验，不只报告最好数字；
  \item \term{优化有边界：}任何性能收益都不能以破坏协议和正确性为代价；
  \item \term{表达面向读者：}先解释问题为什么存在，再介绍实现如何解决，而不是直接堆叠代码细节。
\end{enumerate}

\subsection{本月进度回顾}

\begin{table}[htbp]
  \centering
  \caption{按周回顾}
  \begin{tabularx}{\textwidth}{>{\bfseries}p{0.13\textwidth}X X}
    \toprule
    时间 & 学习/实践重点 & 结果 \\
    \midrule
    第一周 & 昇腾架构、AI Core、GM/UB 与基础 Ascend C & 建立算子执行和数据搬运的基本认识 \\
    第二周 & Ascend C 算子结构、流水线同步、双缓冲和编译流程 & 能够独立阅读、修改并编译算子代码 \\
    第三周 & UDMA/URMA、Ready 协议、Workspace 与多 rank 同步 & 理解 Combine 的跨核、跨卡数据通路 \\
    第四周 & Combine 开发、正确性验证、性能 Profile 与 A/B 对比 & 完成算子实践并形成完整实验闭环 \\
    \bottomrule
  \end{tabularx}
\end{table}

\section{后续计划}

下一阶段计划从“完成一个算子”继续走向“能够稳定地优化和维护算子”：

\begin{enumerate}
  \item 继续补充 Ascend C 指令、内存对齐和流水线调度知识，减少对经验性试验的依赖；
  \item 深入分析 BS128 下 Send 与 Receive 的关键路径，研究 route 粒度流水化和轮询等待比例；
  \item 完善自动化正确性回归，覆盖更多 shape、rank 数和长时间重复运行；
  \item 将实验记录模板化，使源码、编译参数、运行配置和报告可以自动关联；
  \item 学习 Dispatch 与专家计算侧实现，从单个 Combine 算子扩展到完整 MoE 数据通路。
\end{enumerate}

后续优化仍将遵循本月建立的原则：先保证协议正确，再讨论性能；先统一测量边界，再比较
数字；既关注局部阶段，也始终回到多 rank 端到端关键路径。

\section{结语}

本月实习从昇腾架构和 Ascend C 入门，逐步进入片上数据搬运、流水线同步和设备间通信，
最终在 TileXR EP Combine 算子上完成了综合实践。这个过程让我认识到，高性能 AI 算子
并不是一段孤立的数学计算代码，而是硬件存储层次、编译器、通信队列、同步协议和实验
方法共同组成的系统。

从 AI 计算到超节点，核心问题始终是如何让更多计算单元高效、正确地协同工作；从超节点
落到 Combine，工程任务则变成了每一块数据如何被打包、何时发送、如何确认到达、何时
可以复用。通过本月实践，我已经建立起处理这类问题的基本框架，并能够以可验证、可复现
的方式推进后续工作。

\appendix

\section{术语速查}

\begin{longtable}{>{\bfseries}p{0.18\textwidth}p{0.72\textwidth}}
  \toprule
  术语 & 通俗解释 \\
  \midrule
  AI Core & 昇腾芯片中执行 AI 计算和数据搬运的核心单元 \\
  Ascend C & 面向昇腾 AI Core 的算子编程语言与开发方式 \\
  GM & 容量较大的全局存储，可类比为工厂的大仓库 \\
  UB & AI Core 内的 Unified Buffer，可类比为生产线旁的小料架 \\
  MTE & 负责不同存储层次之间数据搬运的硬件流水线 \\
  MoE & 混合专家模型，每个 token 只选择少数专家参与计算 \\
  EP & 专家并行，把不同专家分布在不同设备上 \\
  Dispatch & 把 token 发送到被选中专家的过程 \\
  Combine & 把专家结果送回并按权重聚合的过程 \\
  rank & 分布式任务中的一个参与者，在本报告中通常对应一张卡上的进程 \\
  route & 一条专家输出及其目标 token、目标 rank 等路由信息 \\
  UDMA/URMA & 面向注册内存的设备间直接数据传输能力 \\
  Ready flag & 生产者通知消费者“某一轮数据已经可以使用”的状态标志 \\
  Ping-pong & 两块缓冲区交替使用，以重叠相邻轮次或避免立即覆盖 \\
  Profile & 通过插桩记录各阶段时间和计数，用于定位性能瓶颈 \\
  Strict time & 尽量减少插桩后测得的 Kernel 关键路径时间 \\
  \bottomrule
\end{longtable}

\section{实验复现检查清单}

\begin{enumerate}
  \item 记录源码版本、关键文件哈希和本次 variant 名称；
  \item 记录 CANN、驱动、目标芯片和实际 BiSheng 编译参数；
  \item 记录 rank 数、batch size、hidden size、TopK 和核角色配置；
  \item 编译完成后确认运行加载的是本次生成的 Kernel 和动态库；
  \item 先执行正确性 smoke test，再执行多轮和多 rank 验证；
  \item 性能测试隔离其他任务，记录无效样本和异常原因；
  \item 使用 max-core、max-rank、median-launch 的统一聚合口径；
  \item 同时保留 raw cycles、汇总表、逐 rank 数据和阶段 Profile；
  \item 最终 winner 必须通过双负载验证，不能只依赖单次最好值；
  \item 将未采用方案及原因一并归档，保证结论可以审计。
\end{enumerate}

\begin{thebibliography}{9}
  \bibitem{shazeer2017moe}
  Noam Shazeer et al.
  \newblock Outrageously Large Neural Networks: The Sparsely-Gated Mixture-of-Experts Layer.
  \newblock \emph{International Conference on Learning Representations}, 2017.

  \bibitem{fedus2021switch}
  William Fedus, Barret Zoph, Noam Shazeer.
  \newblock Switch Transformers: Scaling to Trillion Parameter Models with Simple and Efficient Sparsity.
  \newblock \emph{Journal of Machine Learning Research}, 2022.

  \bibitem{ascenddocs}
  华为昇腾社区.
  \newblock Ascend C 算子开发与 CANN 文档中心.
  \newblock \url{https://www.hiascend.com/document}.

  \bibitem{tilexrdesign}
  TileXR 项目.
  \newblock EP URMA Combine 设计说明、测试脚本与性能交付记录，2026.
\end{thebibliography}

\end{document}
